% 对应题号 M2-013，本题条件、变量及几何关系逐项核对。
% 修订：启用中文断行；核对题设及图文对应。
\documentclass[tikz,border=14pt]{standalone}
\usepackage{fontspec,xeCJK}
\usepackage{amsmath,amssymb}
\usepackage{xcolor}
\usepackage{tikz}
\usetikzlibrary{arrows.meta,calc,patterns,positioning,bending}
\IfFontExistsTF{Heiti SC}{\IfFontExistsTF{Heiti SC}{\setmainfont{Heiti SC}}{\setmainfont{FandolHei-Regular}}\IfFontExistsTF{Heiti SC}{\setsansfont{Heiti SC}}{\setsansfont{TeX Gyre Heros}}}{\IfFontExistsTF{Songti SC}{\setmainfont{Songti SC}}{\setmainfont{FandolSong-Regular}}\IfFontExistsTF{Songti SC}{\setsansfont{Songti SC}}{\setsansfont{FandolSong-Regular}}}
\definecolor{bluefield}{RGB}{24,76,255}
\definecolor{redfield}{RGB}{232,49,35}
\definecolor{orangefield}{RGB}{236,132,18}
\definecolor{greenfield}{RGB}{18,132,86}
\definecolor{purplefield}{RGB}{112,72,190}
\definecolor{linegray}{RGB}{60,68,78}
\definecolor{softblue}{RGB}{238,243,255}
\definecolor{softgreen}{RGB}{238,249,244}
\definecolor{softorange}{RGB}{255,247,235}
\tikzset{
  arrow/.style={-{Latex[length=3.1mm,width=2.1mm]},line width=1.1pt},
  thinarr/.style={-{Latex[length=2.5mm,width=1.7mm]},line width=.85pt},
  axis/.style={-{Latex[length=2.5mm,width=1.7mm]},line width=.8pt,draw=linegray},
  guide/.style={draw=linegray!55,dashed,line width=.65pt},
  box/.style={draw=linegray,rounded corners=3pt,fill=white,align=center,inner sep=6pt,line width=.75pt},
  bluebox/.style={box,draw=bluefield,fill=softblue},
  greenbox/.style={box,draw=greenfield,fill=softgreen},
  orangebox/.style={box,draw=orangefield,fill=softorange}
}
\IfFontExistsTF{Heiti SC}{\setCJKmainfont[AutoFakeBold=1.5]{Heiti SC}}{\setCJKmainfont[AutoFakeBold=1.5]{FandolHei-Regular}}
\IfFontExistsTF{Heiti SC}{\setCJKsansfont{Heiti SC}}{\setCJKsansfont{FandolHei-Regular}}
\begin{document}

\begin{tikzpicture}[x=1cm,y=1cm]
\path[use as bounding box] (-10,-4.6) rectangle (10,4.6);
\fill[white] (-10,-4.6) rectangle (10,4.6);
\node[font=\Large\bfseries,text width=18.6cm,align=center] at (0,4.05) {鞍形静电场：场线与等势线构成两族正交双曲线};
\begin{scope}[shift={(-4.6,-.2)}]
\draw[axis] (-3,0)--(3,0) node[right] {$x$}; \draw[axis] (0,-2.7)--(0,2.7) node[above] {$y$};
\foreach \c in {.45,.9,1.4}{\draw[bluefield,line width=.9pt,domain=.5:2.8,samples=60] plot (\x,{\c/\x}); \draw[bluefield,line width=.9pt,domain=-2.8:-.5,samples=60] plot (\x,{\c/\x}); \draw[bluefield,line width=.9pt,domain=.5:2.8,samples=60] plot (\x,{ -\c/\x}); \draw[bluefield,line width=.9pt,domain=-2.8:-.5,samples=60] plot (\x,{ -\c/\x});}
\draw[redfield,line width=1pt] plot[domain=-2.3:2.3,samples=80] (\x,{sqrt(\x*\x+1)}); \draw[redfield,line width=1pt] plot[domain=-2.3:2.3,samples=80] (\x,{-sqrt(\x*\x+1)});
\node[bluefield,fill=white,inner sep=3pt] at (-1.5,3.1) {场线 $xy=C$}; \node[redfield,fill=white,inner sep=3pt] at (0,-3.2) {等势线 $x^2-y^2=C'$};
\foreach\x/\y in{-1.6/-1.2,-1.6/1.2,1.6/-1.2,1.6/1.2}{\draw[arrow,redfield](\x,\y)--({1.18*\x},{.82*\y});}\end{scope}
\begin{scope}[shift={(4.7,-.1)}]
\node[bluebox,text width=6.7cm] at (0,1.65) {$\nabla\cdot\mathbf E=a-a=0$，$\nabla\times\mathbf E=0$};
\node[greenbox,text width=6.7cm] at (0,0) {$\displaystyle \varphi=-\frac a2(x^2-y^2)+C$};
\node[orangebox,text width=6.7cm] at (0,-1.7) {原点处 $\mathbf E=0$，但沿 $x$、$y$ 方向分别表现为不稳定与稳定。};
\end{scope}
\end{tikzpicture}
\end{document}
